\documentclass[11pt,a4paper]{ctexart}
\usepackage[margin=24mm,top=25mm,bottom=24mm]{geometry}
\usepackage{amsmath,amssymb,amsthm,mathtools,bm}
\usepackage{booktabs,tabularx,array}
\usepackage{hyperref,xcolor,microtype}
\usepackage{enumitem,fancyhdr,lastpage,tcolorbox}
\tcbuselibrary{breakable}
\hypersetup{colorlinks=true,urlcolor=blue,linkcolor=blue!50!black}
\setlength{\parindent}{2em}
\setlength{\parskip}{.18em}
\setlist{topsep=.3em,itemsep=.15em}
\definecolor{navy}{HTML}{193149}
\definecolor{teal}{HTML}{1B7281}
\definecolor{paper}{HTML}{F3F7F8}
\pagestyle{fancy}\fancyhf{}
\fancyhead[L]{\small\textcolor{navy}{有限样本认证与最优值收敛}}
\fancyhead[R]{\small\textcolor{navy}{论文 \S7 专项证明}}
\fancyfoot[C]{\small 第 \thepage\ 页 / 共 \pageref{LastPage} 页}
\renewcommand{\headrulewidth}{.3pt}
\newtheorem{theorem}{定理}[section]
\newtheorem{lemma}[theorem]{引理}
\newtheorem{proposition}[theorem]{命题}
\newtheorem{corollary}[theorem]{推论}
\theoremstyle{definition}
\newtheorem{remark}[theorem]{说明}
\newenvironment{proofcn}{\par\noindent\textbf{证明.}\ }{\hfill$\square$\par}
\newtcolorbox{keybox}{breakable,colback=paper,colframe=teal,boxrule=.6pt,arc=1.5mm,left=2.6mm,right=2.6mm,top=2mm,bottom=2mm}
\newcommand{\calP}{\mathcal P}
\newcommand{\calF}{\mathcal F}
\newcommand{\Rsk}{\mathsf R}
\newcommand{\CGraph}{\mathsf C_0}
\newcommand{\Pp}{\mathbb P}
\newcommand{\Ee}{\mathbb E}
\newcommand{\one}{\mathbf 1}
\newcommand{\epsl}{\varepsilon_{\ell}}
\newcommand{\epsu}{\varepsilon_u}
\newcommand{\hC}{\widehat C}
\newcommand{\evec}{\bm e}
\newcommand{\qvec}{\bm q}
\newcommand{\bvec}{\bm b}
\begin{document}
\thispagestyle{empty}
\vspace*{1.1cm}
\noindent{\color{teal}\rule{\textwidth}{1.7pt}}\par
\vspace{.6cm}
{\Huge\bfseries 任务 1：有限样本认证定理（修订版）}\par
\vspace{.24cm}
{\LARGE\bfseries 从无标签分歧率到可认证的稳健加权投票}\par
\vspace{.6cm}
{\Large 论文 \S7 理论正文与证明附录稿}\par
\vspace{.7cm}
\noindent\begin{tabularx}{\textwidth}{@{}p{27mm}X@{}}
\textbf{版本} & 2026 年 10 月 14 日；新增无 big-M 区间 MILP、精确计算认证与可检验 $R_0$ 分支。审查修订：$L_K$ 只对阈值规则取最大值；双路径常数改为 $\max\{L_K,L_{0,K}\}$；新增 $K=5$ 闭式认证（定理 \ref{thm:c248}），$K=5$ 不再需要 MILP 与最优性证书。\\
\textbf{数据} & $n_\ell$ 个每成员带标签验证样本，$n_u$ 个无标签验证样本；分类器在验证前固定。\\
\textbf{定理 A} & 只调用结果 6 图证书 MILP：始终高概率有效，在 $R_0$ 中以 $n^{-1/2}$ 速率逼近总体 minimax 值。\\
\textbf{定理 B} & 使用区间 LP 对偶的无 big-M MILP：所有可行参数上均有 $n^{-1/2}$ 速率，须经有理数或区间算术认证。\\
\end{tabularx}
\vfill
\begin{keybox}
\textbf{必须区分的两个结论：}\par
(1) \emph{有效认证}：不管真实参数是否落在 $R_0$ 中，只要同时置信事件成立，所选规则的真实错误率不超过数据依赖的认证值。\par
(2) \emph{收敛至总体最优}：图证书算法需要 $\sum_iq_i\le1$。对整个原区域 $R$，这个结论不成立；若要求全参数 oracle 一致性，须求解完整的置信可行集 minimax 问题。
\end{keybox}
\vspace{.6cm}
\noindent{\small\color{gray} 本报告给出独立的数学论证；[S11] Q3 在上传的选题文档中仅被作为旧框架结果概述，其原始证明文件未附，故不能声称已经逐行复核旧版 Q3。}
\newpage
\tableofcontents
\newpage

\section{设定与总体 oracle}
固定 $K\ge3$ 个二元分类器。令 $X_i=\one\{f_i(U)\ne Y\}$，并定义
\begin{equation}\label{eq:qD}
q_i=\Pp(X_i=1)<\tfrac12,\qquad
D_{ij}=\Pp(f_i(U)\ne f_j(U))=\Pp(X_i\ne X_j),\qquad
e_{ij}=\Pp(X_iX_j=1).
\end{equation}
二元分类器下，分歧率的恒等式精确成立：
\begin{equation}\label{eq:id}
e_{ij}=\tfrac12(q_i+q_j-D_{ij}).
\end{equation}
对边际向量 $\qvec$ 与成对上界 $\bvec$，定义模糊集
\begin{equation}\label{eq:ambig}
\calP(\qvec,\bvec)=\left\{\pi\ge0:\sum_x\pi_x=1,\ \Ee_\pi X_i=q_i,\ \Ee_\pi X_iX_j\le b_{ij}\ (i<j)\right\}.
\end{equation}
令 $\calF_K$ 为所有确定性自对偶错误规则 $F:\{0,1\}^K\to\{0,1\}$，满足 $F(x)+F(\one-x)=1$。固定规则与 minimax oracle 的错误率分别为
\begin{equation}\label{eq:R}
\Rsk_F(\qvec,\bvec)=\max_{\pi\in\calP(\qvec,\bvec)}\Ee_\pi F(X),\qquad
V(\qvec,\bvec)=\min_{F\in\calF_K}\Rsk_F(\qvec,\bvec).
\end{equation}
默认讨论的总体基准是 $V_0:=V(\qvec,\evec)$，其中 $\evec=(e_{ij})$ 为实际成对矩；由定义，$\calP(\qvec,\evec)$ 非空。可以将真实可信的先验上界 $u_{ij}\ge e_{ij}$ 进一步与统计上界取最小，但这不改变大样本极限 $\evec$。

一个关键事实是：$V_0$ \textbf{并不是}真实联合分布下的贝叶斯最优投票错误率。它是只知道真实 $\qvec$ 与 $\evec$，但不知道更高阶依赖时的\emph{人口总体 minimax} 值。选题文档的区域
\begin{equation}\label{eq:R0}
R_0=\left\{\sum_iq_i\le1\right\}\subseteq
R=\left\{\sum_iq_i-(K-2)q_{\min}\le1\right\}
\end{equation}
是图证书精确值定理的适用范围。

\section{统计假设与同时置信区间}
\paragraph{抽样与训练数据隔离。}
对每个成员有 $n_\ell$ 个来自同一目标部署分布的带标签 i.i.d. 观测，用于计算经验错误率 $\hat q_i$；有 $n_u$ 个目标部署分布上的无标签 i.i.d. 观测，各模型在同一无标签样本上预测，从而计算 $\hat D_{ij}$。不同模型的验证样本可以重合，标签样本和无标签样本也不需要彼此独立；只要求每个经验均值沿样本索引具备 Hoeffding 所需的独立性。所有分类器必须在验证数据使用前固定；若用同一数据拟合或选择分类器，应使用独立校准集或另行给出统一泛化界。

记 $m=\binom K2$、$H=\log(2(K+m)/\delta)$，取
\begin{equation}\label{eq:eps}
\epsl=\sqrt{\frac{H}{2n_\ell}},\qquad\epsu=\sqrt{\frac{H}{2n_u}}.
\end{equation}
用 $[t]_a^b:=\max\{a,\min\{b,t\}\}$ 表示截断，构造
\begin{equation}\label{eq:intervals}
\hat q_i^-=[\hat q_i-\epsl]_0^{1/2},\qquad
\hat q_i^+=[\hat q_i+\epsl]_0^{1/2},\qquad
\hat D_{ij}^-=[\hat D_{ij}-\epsu]_0^1.
\end{equation}
并设
\begin{equation}\label{eq:bhat}
\hat b_{ij}=\min\left\{\hat q_i^+,\hat q_j^+,\left[\frac{\hat q_i^++\hat q_j^+-\hat D_{ij}^-}{2}\right]_0^1\right\}.
\end{equation}
如有确定有效的先验上界 $u_{ij}$，可以把它进一步放进式 \eqref{eq:bhat} 的 $\min$；本文以下的收敛结果仍以真实成对值 $e_{ij}$ 为总体基准。

\begin{lemma}[同时覆盖、偏差大小、保守可行性]\label{lem:coverage}
令事件
\[
\mathcal E_\delta=\left\{\max_i|\hat q_i-q_i|\le\epsl,\ \max_{i<j}|\hat D_{ij}-D_{ij}|\le\epsu\right\}.
\]
则 $\Pp(\mathcal E_\delta)\ge1-\delta$。在该事件上，所有 $i<j$ 同时满足
\begin{align}\label{eq:bounds}
\hat q_i^-\le q_i\le\hat q_i^+,&\qquad0\le \hat q_i^+-q_i\le2\epsl,\\
e_{ij}\le\hat b_{ij},&\qquad 0\le \hat b_{ij}-e_{ij}\le2\epsl+\epsu.\notag
\end{align}
特别地，以下经验置信模糊集必非空，且包含真实联合分布 $\pi^0$：
\begin{equation}\label{eq:Phat}
\widehat{\calP}=\left\{\pi\ge0:\ \sum_x\pi_x=1,\ \hat q_i^-\le\Ee_\pi X_i\le\hat q_i^+,\ \Ee_\pi X_iX_j\le\hat b_{ij}\right\}.
\end{equation}
\end{lemma}
\begin{proofcn}
每个经验变量均在 $[0,1]$ 内，Hoeffding 不等式给出 $\Pp(|\hat q_i-q_i|>\epsl)\le2e^{-2n_\ell\epsl^2}=\delta/(K+m)$，对于分歧率亦然。对 $K+m$ 个变量取并集界得到 $\Pp(\mathcal E_\delta)\ge1-\delta$。

真值 $0<q_i<1/2$，所以截断不破坏 $q_i\in[\hat q_i^-,\hat q_i^+]$；且 $q_i^+-q_i\le2\epsl$。同理 $D_{ij}-2\epsu\le\hat D_{ij}^-\le D_{ij}$。记未截断的上界 $u_{ij}^+=(\hat q_i^++\hat q_j^+-\hat D_{ij}^-)/2$，利用式 \eqref{eq:id} 得
\[
0\le u_{ij}^+-e_{ij}=\tfrac12\left[(\hat q_i^+-q_i)+(\hat q_j^+-q_j)+(D_{ij}-\hat D_{ij}^-)\right]\le2\epsl+\epsu.
\]
因 $e_{ij}\le q_i,q_j\le\hat q_i^+,\hat q_j^+$，做非负截断和取 $\min$ 仍保留 $\hat b_{ij}\ge e_{ij}$，且不增大上偏差。真实 $\pi^0$ 满足所有上述区间与上界，故 $\pi^0\in\widehat{\calP}$。
\end{proofcn}

\section{结果 5：有限片段结构导出的 Lipschitz 常数}
给每条 self-dual 规则 $F$ 写出固定规则的对偶多面体
\begin{equation}\label{eq:dual}
\mathcal D_F=\left\{(a,\beta,\gamma):\ a+\sum_i\beta_i x_i+\sum_{i<j}\gamma_{ij}x_ix_j\ge F(x)\ (\forall x),\ \gamma\ge0\right\}.
\end{equation}
目标 $L_{a,\beta,\gamma}(q,b)=a+\sum_i\beta_iq_i+\sum_{i<j}\gamma_{ij}b_{ij}$。在参数 $(q,b)$ 使 $\calP(q,b)$ 非空时，线性规划强对偶说明
\begin{equation}\label{eq:vertexF}
\Rsk_F(q,b)=\min_{v\in\operatorname{vert}(\mathcal D_F)}L_v(q,b).
\end{equation}
这是上传文档结果 5 的核心内容。$\mathcal D_F$ 与参数 $q,b$ 无关，并且是尖多面体（不存在整条双向直线）：沿任意双向可行方向必须有 $\Delta\gamma=0$，又由全部二元 $x$ 的约束可知 $\Delta a+\sum_i\Delta\beta_i x_i=0$，因而所有方向分量为零。故其顶点有限，有限最优目标必可由顶点取得。

记 $\calF_K^{\rm thr}\subseteq\calF_K$ 为无平局的有符号线性阈值 self-dual 规则（$K=5$ 时共 $1{,}882$ 条）。由选题文档结果 2(b)（阈值充分性），总体 oracle 规则总可以取在 $\calF_K^{\rm thr}$ 中，$\CGraph$ 与 $V$ 中对 $F$ 的最小化也都可以限制在 $\calF_K^{\rm thr}$ 上。令 $d=1+K+m$，定义
\begin{equation}\label{eq:LK}
L_K:=\max_{F\in\calF_K^{\rm thr}}\ \max_{v=(a,\beta,\gamma)\in\operatorname{vert}(\mathcal D_F)}\left(\sum_i|\beta_i|+\sum_{i<j}\gamma_{ij}\right)<\infty.
\end{equation}
又记零截距多面体
\[
\mathcal D_F^0=\{(s,\gamma):\sum_is_ix_i+\sum_{i<j}\gamma_{ij}x_ix_j\ge F(x)\ \forall x,\ \gamma\ge0\},
\]
以及同样定义的 $L_{0,K}$（对 $F\in\calF_K^{\rm thr}$ 的 $\mathcal D_F^0$ 顶点取 $\sum_i|s_i|+\sum\gamma_{ij}$ 的最大值）。令
\begin{equation}\label{eq:C0}
\CGraph(q,b):=\min\{\min_iq_i,B_{\rm graph}(b)\}
=\min_{F\in\calF_K^{\rm thr}}\ \min_{(s,\gamma)\in\mathcal D_F^0}\left(\sum_iq_is_i+\sum_{i<j}b_{ij}\gamma_{ij}\right).
\end{equation}
第二个等号是选题文档的结果 3(a)：\emph{零截距证书最优值恰为独裁或最优图证书}，对所有非负 $(q,b)$ 成立。零截距可行证书必有 $s_i\ge0$，因为把 $x=e_i$ 代入约束。

\begin{lemma}[两种值函数的全局 Lipschitz 性]\label{lem:lip}
在参数对均可行的区域，对固定 $F\in\calF_K^{\rm thr}$ 有
\begin{equation}\label{eq:LipF}
|\Rsk_F(q,b)-\Rsk_F(q',b')|\le L_K\Vert(q,b)-(q',b')\Vert_\infty.
\end{equation}
对任意非负参数（无需对应可行的联合分布），有
\begin{equation}\label{eq:LipC}
|\CGraph(q,b)-\CGraph(q',b')|\le L_{0,K}\Vert(q,b)-(q',b')\Vert_\infty.
\end{equation}
更明确地，上述常数可以保守地选择为 $L_K,L_{0,K}\le (K+m)d^{d/2}$，其中 $d=1+K+m$。
\end{lemma}
\begin{proofcn}
有限个仿射函数 $A_t$ 的最小值在统一斜率上界下是 Lipschitz 的：
\[|\min_tA_t(z)-\min_tA_t(z')|\le\max_t|A_t(z)-A_t(z')|.\]
将其用于式 \eqref{eq:vertexF} 与式 \eqref{eq:C0} 即得两条不等式。为给出不依赖实验枚举的粗显式常数，观察所有顶点都是线性约束矩阵的一组满秩活动行求出的唯一解，矩阵系数属于 $\{-1,0,1\}$，右端属于 $\{0,1\}$。由 Cramer 法则，分母是非零整数，故绝对值至少为 $1$；由 Hadamard 行列式界，每个坐标的分子绝对值至多 $d^{d/2}$。因此各斜率坐标均有 $d^{d/2}$ 的统一上界，乘以 $K+m$ 即得结论。此显式界非常宽松，其用途仅是证明常数有限，而非推荐直接用于样本复杂度估算。
\end{proofcn}


\begin{remark}[$K=5$ 的紧化常数与可复核性]\label{rem:L5}
对 $K=5$，
\begin{equation}\label{eq:L5}
L_5=65,\qquad L_{0,5}=5.
\end{equation}
$L_5$ 是在 $1{,}882$ 条阈值规则（按 $S_5$ 置换共 60 类）对偶多面体的全部 $17{,}508$ 个精确有理顶点上取系数 $\ell_1$ 范数的最大值，由权重为 $(-1,-1,-2,-2,-3)$ 的规则取到。审查已用 cddlib 的 GMP 有理枚举独立复算（选题文档 [S9] \texttt{task1\_review/constants\_K5.py}）。$L_{0,5}=5$ 不需要枚举：由选题文档结果 3(c)，$K=5$ 时 $\CGraph$ 是 5 个独裁项与 22 个奇环代价的最小值，各项系数 $\ell_1$ 范数至多为五环的 $5$。这里的常数都对应 $\ell_\infty$ 参数扰动与系数 $\ell_1$ 范数，不能与其他范数下的常数混用。
\end{remark}

\begin{remark}[与 $K$ 的关系]
下述 $O(\sqrt{\log(K/\delta)/n})$ 结论采用\emph{固定 $K$} 的渐近约定；隐藏常数可以随 $K$ 快速增长。若让 $K=K(n)$ 随样本数增大，必须额外控制 $L_K$，不能据此声称有与 $K$ 无关的统一常数。
\end{remark}

\section{定理 A：使用结果 6 图证书算法的认证与收敛}
定义纯图式数据依赖的认证值
\begin{equation}\label{eq:graphcert}
\hC_G:=\min\left\{\min_i\hat q_i^+,\ B_{\rm graph}(\hat b)\right\}=\CGraph(\hat q^+,\hat b).
\end{equation}
若第一项较小，选择相应成员的独裁投票；否则使用结果 6 的精确 MILP 求得图权重 $y\ge0$，选择按加权度数 $d_i(y)=\sum_{j\ne i}y_{ij}$ 进行的 self-dual 投票。若图证书的最优值以数值精度 $\eta_{\rm alg}$ 近似，则返回经过可行性检查的实际权重及其上界 $\widetilde C_G\in[\hC_G,\hC_G+\eta_{\rm alg}]$。理论上采用精确求解，或者把数值误差作为附加项。

\begin{theorem}[图认证的有限样本覆盖与 oracle 速率]\label{thm:graph}
在本报告抽样假设下，以至少 $1-\delta$ 的概率，上述算法选择的 $\hat F_G$ 满足
\begin{equation}\label{eq:graphcov}
\Pp_{\pi^0}\{\hat F_G(X)=1\}\le\widetilde C_G.
\end{equation}
并且若真实参数位于 $R_0$，即 $\sum_iq_i\le1$，则在同一事件上
\begin{equation}\label{eq:graphgap}
0\le\hC_G-V_0\le L_{0,K}(2\epsl+\epsu),\qquad
0\le\widetilde C_G-V_0\le L_{0,K}(2\epsl+\epsu)+\eta_{\rm alg}.
\end{equation}
因此对于固定 $K$，精确算法满足
\[
\hC_G-V_0=O\!\left(\sqrt{\frac{\log(K/\delta)}{n_\ell}}+\sqrt{\frac{\log(K/\delta)}{n_u}}\right)
\quad\text{（至少以 $1-\delta$ 的概率）}.
\]
\end{theorem}
\begin{proofcn}
在事件 $\mathcal E_\delta$ 上，真实错误率 $q_i\le\hat q_i^+$，因此所选独裁规则的实际风险不超过对应上界。对于任意图权重，图证书上界给出
\[\Rsk(\hat F_G;\pi^0)\le\frac{\sum_{i<j}e_{ij}y_{ij}}{g(y)}\le\frac{\sum_{i<j}\hat b_{ij}y_{ij}}{g(y)}.\]
故经过核验的候选规则，其实际风险均不超过返回的认证值，得到式 \eqref{eq:graphcov}。

由引理 \ref{lem:coverage}，$\hat q_i^+\ge q_i$，$\hat b_{ij}\ge e_{ij}$；又因 $B_{\rm graph}$ 对每个 $b_{ij}$ 均单调不减，得 $\hC_G\ge\CGraph(q,e)$。真实参数位于 $R_0$ 时，选题文档结果 3(b) 说明 $\CGraph(q,e)=V_0$。另一方面，由式 \eqref{eq:bounds}，
\[
\Vert(\hat q^+,\hat b)-(q,e)\Vert_\infty\le2\epsl+\epsu.
\]
对式 \eqref{eq:C0} 应用引理 \ref{lem:lip} 得 $\hC_G\le V_0+L_{0,K}(2\epsl+\epsu)$。算法误差至多再增加 $\eta_{\rm alg}$。
\end{proofcn}

\begin{remark}[为什么不需要额外进行模型选择修正？]
证书同时覆盖所有成员与所有模型对；在覆盖事件上，\emph{每一个}合法图权重都满足认证不等式。因此 $y$ 虽然用同一批验证数据自适应挑选，也不需要再对候选图或候选规则数做并集界。这一点与“先选一个规则再用同一数据做未经调整的经验风险区间”不同。
\end{remark}


\subsection{可检验的 $R_0$ 条件及自动分支}
定义事件 $\mathcal T_0:=\{\sum_i\hat q_i^+\le1\}$。在同时覆盖事件 $\mathcal E_\delta$ 上，
\begin{equation}\label{eq:testR0}
\mathcal T_0\ \Longrightarrow\ \sum_iq_i\le\sum_i\hat q_i^+\le1.
\end{equation}
因此\textbf{数据满足 $\mathcal T_0$ 时}，允许采用结果 6 的图证书 MILP，并且这时取得总体 minimax 的 $O(n^{-1/2})$ 保证。若 $\mathcal T_0$ 不满足，则不能反推真实参数不属于 $R_0$，只能退回完整区间求解。由于两条路径共用同一个同时置信覆盖事件，不需要额外为数据驱动分支扣除一份置信概率。

\section{定理 B：完整置信模糊集的全参数认证}
若希望\emph{不依赖} $R_0$ 条件，则使用式 \eqref{eq:Phat} 的整个置信可行集，并求解
\begin{equation}\label{eq:Cfull}
\hC_{\rm full}=\min_{F\in\calF_K}\ \max_{\pi\in\widehat{\calP}}\Ee_\pi F(X),
\qquad \hat F_{\rm full}\in\operatorname*{argmin}_{F\in\calF_K}\max_{\pi\in\widehat{\calP}}\Ee_\pi F(X).
\end{equation}
如果经验约束在低概率的置信事件外碰巧不可行，则定义认证值为 $1$ 并任取一条 self-dual 规则作为回退；在覆盖事件上由引理 \ref{lem:coverage} 可行性自动保证。该问题仍是有限维的精确优化，内层 LP 中含有 $q_i^-$ 和 $q_i^+$ 两组边际不等式。对于固定 $F$，对偶具有明确形式：
\begin{align}\label{eq:intervaldual}
\widehat{\mathsf R}_F=\min_{a,\lambda^+,\lambda^-,\gamma}\quad& a+\sum_i(\lambda_i^+\hat q_i^+-\lambda_i^-\hat q_i^-)+\sum_{i<j}\gamma_{ij}\hat b_{ij}\\
\text{s.t.}\quad& a+\sum_i(\lambda_i^+-\lambda_i^-)x_i+\sum_{i<j}\gamma_{ij}x_ix_j\ge F(x)\quad(\forall x),\notag\\
&\lambda_i^+,\lambda_i^-,\gamma_{ij}\ge0.\notag
\end{align}
所以仍可以比较二次证书在 $x$ 与 $\one-x$ 的差，交叉项相消，推出存在最优的\emph{有符号线性阈值规则}；计算上可枚举 self-dual 规则（或只枚举可实现的阈值规则）并逐一解 LP。\textbf{结果 6 的纯图 MILP 不能替代}式 \eqref{eq:Cfull}。


\subsection{结果 2(a) 的无 big-$M$ 区间 MILP}\label{subsec:milp}
将区间边际的自由对偶系数写成 $\beta_i=\lambda_i^+-\lambda_i^-$，其中 $\lambda_i^\pm\ge0$。令 $\mathcal Z$ 为所有 $\{x,\one-x\}$ 互补模式对中任取一个代表的集合，$|\mathcal Z|=2^{K-1}$。对每个 $z\in\mathcal Z$ 引入一个二元变量 $t_z$：$t_z=1$ 表示 $F(z)=1$，而 $F(\one-z)=0$。定义
\[
Q(x)=a+\sum_i(\lambda_i^+-\lambda_i^-)x_i+\sum_{i<j}\gamma_{ij}x_ix_j.
\]
\begin{proposition}[完整区间最优值的无 big-$M$ 精确混合整数线性规划]\label{prop:milp}
若 $\widehat{\calP}\ne\varnothing$，则 $\hat C_{\rm full}$ 精确等于
\begin{equation}\label{eq:milp}
\begin{aligned}
\min_{a,\lambda^+,\lambda^-,\gamma,t}\quad& a+\sum_i\bigl(\lambda_i^+\hat q_i^+-\lambda_i^-\hat q_i^-\bigr)+\sum_{i<j}\gamma_{ij}\hat b_{ij},\\
\text{s.t.}\quad& Q(z)\ge t_z,\qquad\ \forall z\in\mathcal Z,\\
&Q(\one-z)\ge1-t_z,\quad\forall z\in\mathcal Z,\\
& t_z\in\{0,1\},\quad \lambda_i^\pm,\gamma_{ij}\ge0,\quad a\in\mathbb R.
\end{aligned}
\end{equation}
这里\emph{没有任何 big-$M$ 常数}，且每对互补模式只需一个二元变量。对于 $K=5$，共有 $16$ 个二元变量、$21$ 个连续变量以及 $32$ 条互补模式约束（不计变量界）。
\end{proposition}
\begin{proofcn}
对任意固定 self-dual $F$，令 $t_z=F(z)$，则两条互补不等式正好是式 \eqref{eq:intervaldual} 中该互补对上的 $Q\ge F$ 条件；故每个 $F$ 的 LP 对偶是本 MILP 固定 $t$ 后的连续部分。反向地，任意二元向量 $t$ 直接定义一个 self-dual $F$。对 $t$ 及连续对偶变量同时取最小，相当于对全部规则的最坏风险取最小。因为 $\widehat{\calP}$ 非空，固定规则满足 LP 强对偶，从而两边相等。
\end{proofcn}

式 \eqref{eq:intervaldual} 的 $Q(x)-Q(\one-x)$ 仍是
\[
\sum_i \left(\lambda_i^+-\lambda_i^-+\tfrac12\sum_{j\ne i}\gamma_{ij}\right)(2x_i-1),
\]
因而仍存在最优的\emph{有符号齐次线性阈值规则}。$K=5$ 时可预先枚举 $1{,}882$ 条不同的无平局 self-dual 阈值规则；逐条求解式 \eqref{eq:intervaldual} 的 LP，最后取最小值，作为 MILP 的独立交叉核验。前期提供的计数为 $1{,}882$；此处不把代码枚举视作已经完成的独立实验。

\begin{theorem}[完整区间鲁棒方法：所有参数的高概率 oracle 界]\label{thm:full}
在同样抽样假设下，以至少 $1-\delta$ 的概率，\emph{不要求}真实参数处于 $R$ 或 $R_0$，有
\begin{equation}\label{eq:fullcov}
\Pp_{\pi^0}\{\hat F_{\rm full}(X)=1\}\le\hC_{\rm full},
\end{equation}
并且
\begin{equation}\label{eq:fullgap}
\boxed{\quad 0\le\hC_{\rm full}-V_0\le L_K(2\epsl+\epsu).\quad}
\end{equation}
因此对固定 $K$，$\hC_{\rm full}\to V_0$ 的速率同样为 $n_\ell^{-1/2}+n_u^{-1/2}$（带 $\sqrt{\log(K/\delta)}$ 系数）。
\end{theorem}
\begin{proofcn}
考虑覆盖事件 $\mathcal E_\delta$，引理 \ref{lem:coverage} 给出 $\pi^0\in\widehat{\calP}$，所以对\emph{任意}可能依赖验证数据的 $F$，都有 $\Ee_{\pi^0}F\le\max_{\pi\in\widehat{\calP}}\Ee_\pi F$；对 $\hat F_{\rm full}$ 应用得到式 \eqref{eq:fullcov}。

同样由于总体真值约束集 $\calP(q,e)\subseteq\widehat{\calP}$，可对全部 $F$ 取最小，得到 $\hC_{\rm full}\ge V_0$。取阈值形式的总体 oracle 规则 $F^*\in\arg\min_{F\in\calF_K^{\rm thr}}\Rsk_F(q,e)$（由阈值充分性，它也是全部 $\calF_K$ 上的最优规则），则
\begin{align*}
\hC_{\rm full}&\le\max_{\pi\in\widehat{\calP}}\Ee_\pi F^*(X)\\
&=\max_{q'\in\mathcal I:\ \calP(q',\hat b)\ne\varnothing}\Rsk_{F^*}(q',\hat b),
\end{align*}
其中 $\mathcal I=\prod_i[\hat q_i^-,\hat q_i^+]$，等式把经验可行集按边际 $q'$ 分解。
对任何出现在最大化中的 $q'$，都满足 $\Vert q'-q\Vert_\infty\le2\epsl$，而 $\Vert\hat b-e\Vert_\infty\le2\epsl+\epsu$。由于 $\calP(q',\hat b)$ 和 $\calP(q,e)$ 都非空，应用引理 \ref{lem:lip} 中固定 $F^*$ 的全局 Lipschitz 性，
\[
\Rsk_{F^*}(q',\hat b)\le\Rsk_{F^*}(q,e)+L_K(2\epsl+\epsu)
=V_0+L_K(2\epsl+\epsu).
\]
对 $q'$ 取最大，即得式 \eqref{eq:fullgap}。
\end{proofcn}

\begin{corollary}[数据选择规则的 oracle 逼近]\label{cor:sel}
在定理 \ref{thm:graph} 的 $R_0$ 条件下，$\hat F_G$ 满足
$\Ee_{\pi^0}\hat F_G\le V_0+L_{0,K}(2\epsl+\epsu)+\eta_{\rm alg}$。
对任何可行真实参数，$\hat F_{\rm full}$ 满足
$\Ee_{\pi^0}\hat F_{\rm full}\le V_0+L_K(2\epsl+\epsu)$。
\end{corollary}
\begin{proofcn}
分别合并两条定理的覆盖保证与认证值上界。
\end{proofcn}


\subsection{$K=5$ 的闭式认证：不需要 MILP}\label{subsec:c248}
选题文档结果 8 给出：对一切可行参数，$V(q,b)=\min_{j\le248}L_{Q_j}(q,b)$，其中 $Q_j=(a_j,\beta_j,\gamma_j)$ 是 248 个显式二次证书，分属 12 个 $S_5$ 轨道。每个 $Q_j$ 对某条阈值规则 $F_j$（权重见附带的 \texttt{U248\_certificates.json}）满足 $Q_j\ge F_j$ 与 $\gamma_j\ge0$。定义
\begin{equation}\label{eq:c248}
\hC_{248}:=\min_{j\le248}\Bigl[a_j+\sum_i\bigl(\beta_{ij}^+\hat q_i^+-\beta_{ij}^-\hat q_i^-\bigr)+\sum_{i<k}\gamma_{ik,j}\hat b_{ik}\Bigr],
\end{equation}
其中 $t^\pm=\max\{\pm t,0\}$。选出取到最小值的 $j^\dagger$，并采用规则 $\hat F_{248}:=F_{j^\dagger}$。

\begin{theorem}[$K=5$ 的闭式有限样本认证]\label{thm:c248}
在本报告的抽样假设下，以至少 $1-\delta$ 的概率，对一切可行的总体参数（不要求 $R$ 或 $R_0$），
\[
\Pp_{\pi^0}\{\hat F_{248}(X)=1\}\le\hC_{248},\qquad
\hC_{\rm full}\le\hC_{248},\qquad
0\le\hC_{248}-V_0\le 47\,(2\epsl+\epsu).
\]
其中 $47=\max_j\bigl(\sum_i|\beta_{ij}|+\sum_{i<k}\gamma_{ik,j}\bigr)$。
\end{theorem}
\begin{proofcn}
取 $\lambda^+=\beta_j^+$、$\lambda^-=\beta_j^-$，则 $(a_j,\lambda^+,\lambda^-,\gamma_j)$ 是规则 $F_j$ 的区间对偶式 \eqref{eq:intervaldual} 的可行点，目标值正是式 \eqref{eq:c248} 中的第 $j$ 项。所以每一项都不小于 $\max_{\pi\in\widehat\calP}\Ee_\pi F_j$。由此得到 $\hC_{\rm full}\le\hC_{248}$；在覆盖事件上 $\pi^0\in\widehat\calP$，因此得到第一个不等式。下界 $\hC_{248}\ge\hC_{\rm full}\ge V_0$ 同定理 \ref{thm:full}。

对上界，取 $j^*$ 使 $L_{Q_{j^*}}(q,e)=V_0$（由结果 8，在全部可行参数上成立）。在覆盖事件上，$0\le\hat q_i^+-q_i\le2\epsl$，$0\le q_i-\hat q_i^-\le2\epsl$，$0\le\hat b_{ik}-e_{ik}\le2\epsl+\epsu$。于是第 $j^*$ 项减去 $L_{Q_{j^*}}(q,e)$ 等于
\[
\sum_i\bigl[\beta_{ij^*}^+(\hat q_i^+-q_i)+\beta_{ij^*}^-(q_i-\hat q_i^-)\bigr]+\sum_{i<k}\gamma_{ik,j^*}(\hat b_{ik}-e_{ik})\le 47\,(2\epsl+\epsu).
\]
\end{proofcn}

\begin{remark}[计算与严格性]
式 \eqref{eq:c248} 只需对 248 个仿射函数逐一求值，可以全程用有理数完成；附带脚本 \texttt{task1\_C248\_closed\_form.py} 还会对 248 个证书逐状态精确复核 $Q_j\ge F_j$。因此 $K=5$ 时既不需要求解 MILP，也不需要 \S\ref{subsec:rigor} 中的全局最优性下界：定理 \ref{thm:c248} 的速率常数 47 是无条件的，还小于 $L_5=65$。$\hC_{248}$ 与 $\hC_{\rm full}$ 一般不保证相等，因为区间对偶的最优点可能落在某个 $\beta_i=0$ 的非顶点处。不过在 50 个随机有理区间实例上，两者数值相等（选题文档 [S9] \texttt{task1\_review/check\_task1v2.py}）。式 \eqref{eq:milp} 的 MILP 留给一般的 $K$。
\end{remark}

\section{可交付的有限样本自动分支定理}
\begin{theorem}[数据驱动双路径：覆盖、收敛及 $K=5$ 常数]\label{thm:hybrid}
在抽样和固定模型假设下，定义算法：若 $\sum_i\hat q_i^+\le1$，选择经证书核验的图投票，返回上界 $\widetilde C_G$；否则选择经证书核验的式 \eqref{eq:milp} 最优规则，返回上界 $\widetilde C_{\rm full}$。记 $\eta\ge0$ 为经\emph{严格核验}的加性最优性误差，即返回上界与对应精确算法值之差的上界。一般 $K$ 取命题 \ref{prop:lower} 的 $\eta=U-L$；$K=5$ 时两个分支都是闭式，$\eta=0$。记输出为 $(\widehat F,\widetilde C)$，$r=2\epsl+\epsu$。则以至少 $1-\delta$ 的概率，
\begin{equation}\label{eq:hybrid}
\Pp_{\pi^0}\{\widehat F\text{ 出错}\}\le\widetilde C,
\qquad0\le\widetilde C-V_0\le \max\{L_K,L_{0,K}\}\, r+\eta.
\end{equation}
在进入图分支时，右侧可以加强为 $L_{0,K}r+\eta$。对 $K=5$，图分支用结果 3(c) 的 27 项闭式计算 $\hC_G$，完整分支改用定理 \ref{thm:c248} 的 $\hC_{248}$。两个分支都是有限个仿射函数的有理求值，因此 $\eta=0$，并且
\[
\hC_G-V_0\le 5r,\qquad
\hC_{248}-V_0\le47r.
\]
这里的 $\eta$ 不能仅凭求解器的浮点容差声称为严格上界；必须由命题 \ref{prop:lower} 的有理分布 $\pi$ 与有理对偶上界共同给出。
\end{theorem}
\begin{proofcn}
在事件 $\mathcal E_\delta$ 上，若进入图分支，则由式 \eqref{eq:testR0} 确定真实参数位于 $R_0$，从而定理 \ref{thm:graph} 既保证覆盖又保证图证书与 $V_0$ 的差距。若进入完整区间分支，则由定理 \ref{thm:full} 得到不依赖 $R_0$ 的结论。两个分支在同一覆盖事件上逐一成立，所以其数据依赖选择也成立；取两个 Lipschitz 常数中较大的 $\max\{L_K,L_{0,K}\}$ 得到统一式 \eqref{eq:hybrid}。$K=5$ 的两条不等式分别来自定理 \ref{thm:graph}（取 $L_{0,5}=5$，且 $\hC_G$ 为闭式，$\eta=0$）与定理 \ref{thm:c248}。计算加性误差 $\eta$ 只在已经核验优化下界和候选上界的情况下加入。
\end{proofcn}

\section{必要限制、计算与实践细节}
\subsection{图证书不可能在整个 $R$ 内实现人口最优一致性}
第九版选题文档结果 4(a) 的解析反例：$K=5$、$q_i=0.3$、$e_{ij}=0.05$，属于 $R$ 但不属于 $R_0$。此时 $V(q,e)=0$（五人多数保证零错误），而 $B_{\rm graph}(e)>0$（按文档结果 6 的完全对称计算可得 $3e=0.15$）。即便令 $n_\ell,n_u\to\infty$，图证书算法的认证值也收敛到 $0.15$ 而不是 $0$。因此若仅使用结果 6，论文 \S7 \textbf{不得}声称在整个 $R$ 内对 $V(q,e)$ 的 $n^{-1/2}$ 收敛。

\subsection{带先验共错上界的情况}
如果外部先验给出 $u_{ij}\ge e_{ij}$，可定义
$\hat b_{ij}=\min\{u_{ij},\hat q_i^+,\hat q_j^+,(\hat q_i^++\hat q_j^+-\hat D_{ij}^-)/2\}$，覆盖率证明不变。因为即便先验很松，统计项仍收敛到 $e_{ij}$，总体目标仍是 $V(q,e)$。如果\emph{只提供固定先验上界而没有分歧估计}，可认证但通常不能收敛到以精确 $e$ 为参数的 oracle；必须改为以真实先验上界为目标，或者明确不可识别的剩余差距。

\subsection{严格可验证计算：候选上证书与优化下界必须区分}\label{subsec:rigor}
若验证数据的经验频率 $\hat q_i,\hat D_{ij}$ 是有理数，但 Hoeffding 半径包含对数与平方根，推荐对半径\emph{向上舍入为有理数}，并以有理数构造 $\hat q^\pm,\hat D^-,\hat b$。例如取网格步长 $h>0$，置 $\tilde\varepsilon= h\lceil \varepsilon/h\rceil$；将 $h$ 选为 $o(n^{-1/2})$ 以保持速率。严格实数区间算术也可实现同样效果。

对求解器返回的候选 $t$（对应规则 $F$）、有理数化的 $(a,\lambda^+,\lambda^-,\gamma)$，需要对全部 $2^K$ 个输入模式\emph{逐点精确核对} $Q(x)\ge F(x)$，检查非负性，并按有理数计算
\[
U=a+\sum_i(\lambda_i^+\hat q_i^+-\lambda_i^-\hat q_i^-)+\sum_{i<j}\gamma_{ij}\hat b_{ij}.
\]
一旦通过核验，$U$ 就是该规则对整个 $\widehat{\calP}$ 的有效\textbf{上界}，因而足以用于真实风险认证——即便并未证明这条规则是最优规则。\emph{仅将浮点数四舍五入或用 \texttt{limit\_denominator} 拟合为有理数，并不能自动保证对偶可行；必须在转换之后重新精确验算全部约束。}

若还要声称 $U\le\hat C_{\rm full}+\eta$，并据此写出 oracle 收敛速率中的算法项 $\eta$，就需要一个有效的\textbf{优化下界} $L\le\hat C_{\rm full}$。本文不做一般 $K$ 的全局最优性证明，而用下面这个只需一次线性规划的精确下界。

\begin{proposition}[随机化松弛给出的精确下界]\label{prop:lower}
对任意 $\pi\in\widehat{\calP}$，记 $\mathcal Z$ 为互补模式对的代表集，令
\[
L(\pi):=\sum_{z\in\mathcal Z}\min\{\pi_z,\pi_{\one-z}\}.
\]
则 $\hat C_{\rm full}\ge L(\pi)$。$L(\pi)$ 在 $\widehat\calP$ 上的最大值等于随机化 self-dual 规则的区间 minimax 值，也等于式 \eqref{eq:milp} 把 $t_z$ 放松到 $[0,1]$ 后的线性规划值。
\end{proposition}
\begin{proofcn}
对任一确定性 self-dual 规则 $F$，每个互补对中恰有一个模式取 $F=1$，所以 $\Ee_\pi F\ge\sum_z\min\{\pi_z,\pi_{\one-z}\}$。于是
\[
\hat C_{\rm full}=\min_F\max_{\pi'\in\widehat\calP}\Ee_{\pi'}F\ge\min_F\Ee_\pi F\ge L(\pi).
\]
把 $t_z$ 放松到 $[0,1]$ 后，两条互补约束相加得 $Q(z)+Q(\one-z)\ge1$，这正是随机化规则的对偶约束；再由线性规划强对偶，即得后两个等式。
\end{proofcn}

因此，计算时只需给出两个精确的有理对象：(i) 一个经逐状态核验的候选规则及其对偶上界 $U$；(ii) 一个经精确核验属于 $\widehat\calP$ 的有理分布 $\pi$，以及 $L=L(\pi)$。于是 $0\le U-\hat C_{\rm full}\le\eta:=U-L$ 严格成立，定理 \ref{thm:hybrid} 中的 $\eta$ 就取这个值。

附带脚本中的 \texttt{certified\_lower\_bound} 先用对偶单纯形求解随机化线性规划，再由活动约束与正变量用有理数 Gauss--Jordan 消元精确重建最优顶点，最后逐条精确检查 $\pi\in\widehat\calP$。重建失败时只返回平凡下界 $0$，绝不输出未经核验的数值。\texttt{certified\_gap} 同时返回 $U$、$L$ 与 $\eta$。

$\eta$ 恰好是确定性与随机化 minimax 值之差的一个上界，一般不为零。例如五个独立、错误率 $1/5$ 的成员（区间退化为点、$\hat b_{ij}=1/25$）时，$U=3/25$，$L=1/10$，$\eta=1/50$。审查在 $K=5,6,7$ 的随机有理区间实例上做了测试（脚本见选题文档 [S9] 的 \texttt{task1\_review/} 目录）：精确重建全部成功，绝大多数实例上 $\eta=0$，此时候选规则同时被证明是全局最优的。

在置信事件外，若 $\widehat{\calP}$ 为空或求解器无法提供可验证的候选上界，应返回平凡认证值 $1$ 并将状态标记为退回处理，绝不能输出未经验证的数值声称覆盖。

\subsection{可选收紧：经验 Bernstein 与 Clopper--Pearson}
任意对 $q_i$ 与 $D_{ij}$ 有\emph{同时覆盖}的区间族都可替换 Hoeffding 置信区间，定理中仅需把 $r$ 换为实际区间端点对真值的统一上界。具体实现可用各 Bernoulli 样本的 Clopper--Pearson 单侧精确区间（按 $K+\binom K2$ 个指标分配显著性预算），或带适当有限样本修正的经验 Bernstein 上下界。若还希望保留文中\emph{显式}的 $O(n^{-1/2})$ 速率，可同时构造两套覆盖区间，分别分配 $\delta/2$，并把较紧的区间与 Hoeffding 区间相交；在两套区间都覆盖的事件上，交集仍包含真实参数，且区间宽度不超过按 $\delta/2$ 构造的 Hoeffding 区间宽度（半径中的对数项由 $\log(2(K+m)/\delta)$ 变为 $\log(4(K+m)/\delta)$，速率不变）。这种方法通常提升低方差事件的经验紧度，但不能不经计算就保证每一组数据的区间严格缩小。

\subsection{样本数不同、复用数据以及总体分布偏移}
如果各成员有不同的 $n_{\ell,i}$，直接把 $\epsl$ 替换为 $\varepsilon_{\ell,i}=\sqrt{H/(2n_{\ell,i})}$；若各对分歧有不同 $n_{u,ij}$，对应替换为 $\varepsilon_{u,ij}$。最终误差可用这些半径的最大值控制。数据自适应选择\emph{投票规则}是允许的，但不能不加修正地复用同一验证集去\emph{训练基础分类器}。若训练/验证和部署分布不同，本文所有覆盖及一致性结论都需要额外分布偏移条件。

\section{面向论文 \S7 的精简表述}
\begin{keybox}
\textbf{Theorem (Finite-sample certifiable voting).}
Let $K$ fixed binary voters be evaluated on $n_\ell$ i.i.d. labeled and $n_u$ i.i.d. unlabeled samples from a common target distribution. Let $\hat q_i^\pm,\hat D_{ij}^-$ be simultaneous Hoeffding bounds at level $1-\delta$ and put $\hat b_{ij}=\min\{\hat q_i^+,\hat q_j^+,(\hat q_i^++\hat q_j^+-\hat D_{ij}^-)/2\}$. Write $V_0=V(q,e)$ for the population deterministic self-dual minimax error. Then: 
\begin{enumerate}[leftmargin=1.5em]
\item The graph-certificate rule selected by minimizing $\min\{\min_i\hat q_i^+,B_{\rm graph}(\hat b)\}$ has true error at most its reported certificate with probability $\ge1-\delta$. If $\sum_iq_i\le1$, its certificate satisfies $0\le\hat C_G-V_0\le L_{0,K}(2\epsl+\epsu)$.
\item The robust rule selected from $\widehat{\calP}$ has true error at most $\hat C_{\rm full}$ with probability $\ge1-\delta$ for every feasible population parameter, and satisfies $0\le\hat C_{\rm full}-V_0\le L_K(2\epsl+\epsu)$. For $K=5$ the closed-form certificate $\hat C_{248}$ (a minimum of 248 explicit affine functions of $\hat q^\pm,\hat b$) has the same coverage guarantee and $0\le\hat C_{248}-V_0\le 47(2\epsl+\epsu)$, with no optimization required.
\end{enumerate}
Here $L_{0,K},L_K$ (maxima over signed threshold rules; $L_{0,5}=5$, $L_5=65$) are finite constants depending only on $K$, justified by the finite dual-vertex representation, with $2\epsl+\epsu=O(\sqrt{\log(K/\delta)/n_\ell}+\sqrt{\log(K/\delta)/n_u})$ for fixed $K$.
\end{keybox}
\vspace{.4em}
\noindent\textbf{结语：} 任务 1 的统计定理已形成可审查的证明：区间无 big-$M$ MILP 与可检验的图/完整分支都明确。与旧框架不同，成对信息在经验层面只有\emph{上界}；包含真实联合分布的区间模糊集保证了认证覆盖，而不经修改的图算法只在 $R_0$ 具有总体 minimax 一致性。全参数一致性需要完整鲁棒优化；作为认证算法，还须附准确的算术可行性核验，才可正式声明所报告数值的严格有效性。$K=5$ 时，完整鲁棒优化可以换成闭式 $\hC_{248}$：它只需有理数求值，常数为 47。
\end{document}
